% !PRINTED 486
\weberpage{486}
%<<CONT
an algebraic equation exists, and these equations are set up there
for $m=7,\ 23,\ 31,\ 47,\ 71$.

If we put in them:
\[
\omega=\frac{-1}{\sqrt{-m}},\quad m\omega=\sqrt{-m},\quad f(\sqrt{-m})=\sqrt2\,x,
\]
then by \S\,34 we obtain:
\[
A=x^2+\frac{(-1)^{\frac{m+1}{8}}}{x},\quad B=4\,x+\frac{(-1)^{\frac{m+1}{8}}}{x}.
\]
From this there results, for example, for $m=23$, where $A=1$, the equation:
\weq{1}{f(\sqrt{-23})=\sqrt2\,x,\quad x^3-x-1=0}
and for $m=31$:
\[
x^9-4\,x^6+3\,x^3-1=0.
\]
This is in the first place a cubic equation for $x^3$; but from it the
cubic equation for $x$ itself is easily split off:
\weq{2}{f(\sqrt{-31})=\sqrt2\,x,\quad x^3-x^2-1=0.}
In like manner the equations result:
\weq{3}{f(\sqrt{-47})=\sqrt2\,x,\quad x^5-x^3-2\,x^2-2\,x-1=0,}
\weq{4}{f(\sqrt{-71})=\sqrt2\,x,\quad x^7-2\,x^6-x^5+x^4+x^3+x^2-x-1=0.}
In the last case, $m=71$, the factor $(x+1)^2$, which is foreign to the
problem, has been separated off from the equation of degree 9 that is
immediately obtained.

2. In order to arrive at a further mode of calculation, we apply the
transformation of the second order. We have first of all in general
[for a variable $\omega$, \S\,34, (17)]:
\[
\begin{aligned}
f(2\omega)^8+f_2(2\omega)^8&=2\,\frac{f(\omega)^8+f_1(\omega)^8}{f_2(\omega)^4},\\
f(2\omega)^8-f_2(2\omega)^8&=\frac{16}{f_2(\omega)^8},
\end{aligned}
\]
from which:
\[
\begin{aligned}
f(2\omega)^8&=\frac{f(\omega)^8+f_1(\omega)^8}{f_2(\omega)^4}+\frac{8}{f_2(\omega)^8},\\
f_2(2\omega)^8&=\frac{f(\omega)^8+f_1(\omega)^8}{f_2(\omega)^4}-\frac{8}{f_2(\omega)^8};
\end{aligned}
\]
%CONT>>
